\chapter{The Loop That Fails to Close}
% OUTLINE Ch6 "The Loop That Fails to Close (curvature)". Laws: CITATION
% (MTW loop pedagogy; Riemann 1854 lecture + 1861 Commentatio, both
% posthumous; commutator form via Ricci & Levi-Civita 1901; contractions
% w/ Ricci's name), LENS: THE SECOND AND FINAL EXHIBIT rides here under
% Beastmaster's two pre-issued fences (verb = illustrates, never is; the
% GR-name fence = shape only, "does not compute curvature," the face
% still owed) -- if both could not hold, the instruction was drop to
% zero; they hold. Dose closes at 2. No metric anywhere: curvature from
% the connection alone = the no-ruler summit. Gate: Kodo per-claim;
% Beastmaster reads the exhibit sentence on arrival.

Now run the experiment the last chapter could not resist proposing. The
pedagogy is the great gravitation textbook's own --- Misner, Thorne, and
Wheeler teach curvature by exactly this picture --- and it goes like a
bench procedure. Mark a small closed circuit on the manifold: out along
one span, across along a second, back along the first, home along the
second, a little parallelogram of edges. Take an arrow, and carry it
around the circuit by the chosen rule of Chapter 5, faithfully, never
letting it turn relative to the rule's own notion of parallel. The arrow
comes home. Compare it with the arrow that stayed.

On the schoolroom's paper the two agree, always, and that agreement is
what the word flat will now be earned to mean. On the Earth of the last
chapter's experiment they do not agree: the traveler returns rotated,
and --- this is the fact that turns a curiosity into an instrument ---
the disagreement is lawful. Shrink the circuit and the turning shrinks
in proportion to the little parallelogram's area; double the carried
arrow and the difference doubles; swap the circuit's two spans ---
traverse it the other way around --- and the turning reverses sign.
Linear in the arrow carried, linear in each span, antisymmetric under
their exchange: the reader has spent two chapters learning what answers
to that description. It is a machine. There is a tensor whose slots take
the carried arrow and the loop's two spans, and whose answer is the
change the loop inflicts; it is called the Riemann curvature tensor, and
it is the object this construction has been walking toward since the
first chart was laid down.

The trade reaches the same machine by a second route, and the two routes
teach different things, so the book states both. Instead of one loop,
take the two probes separately: differentiate a field of arrows east
then north, by the honest derivative of Chapter 5, and again north then
east, and subtract. On flat paper the order of probing never matters; on
curved ground it does, and the failure of the two second derivatives to
commute is, component for component, the same tensor the loop measures.
That is the form the absolute differential calculus of Ricci and
Levi-Civita computes with --- curvature as the commutator of covariant
derivatives --- and it makes plain a fact the loop picture hides: no
metric was consulted. The carrying rule alone --- the declared structure
of Chapter 5, no ruler anywhere --- suffices to define curvature,
measure it, and compute with it. The longest-running theme of this book
reaches its summit here: the celebrated tensor of curved geometry stands
complete on a manifold that still does not know what a length is. And
the word flat now has its earned meaning: a carrying rule is flat when
every loop returns every arrow as itself --- when the machine reads zero
on all inputs everywhere.

The history is Riemann's, and it is the second half of a story this book
already began. The lecture of 1854 posed the question in nearly this
form: if space is an extended magnitude charted region by region, what
measures its deviation from flatness? The apparatus that answers ---
the four-index expression the trade now writes --- appears in his essay
of 1861 for the Paris Academy, on an applied problem of heat conduction,
where the geometry enters as the tool; and like the lecture, the essay
was published only after his death, in the collected works his friends
assembled. Riemann held the whole subject --- the charts, the question,
and the answering machine --- and lived to publish almost none of it.
The reader who wants the machine's finer anatomy --- its symmetries,
which cut its many components down to few, and its standard contractions,
the smaller readings that carry Ricci's name and, contracted once more,
yield a single scalar --- will find them in the standing references;
this book takes from the anatomy only what later chapters spend.

\begin{fence}
An illustration, not an identity --- and the last device appearance in
this book. The apparatus documented in Volume 1, built to describe its
own act of describing and to close that circuit on itself, returns from
the closed circuit changed: a residue survives the loop, reproducibly,
refusing to vanish however the descent is refined. That is the same
\emph{shape} the reader has just built --- a carrying around a closed
loop that fails to return its object as itself --- and it is offered
here as an illustration of the shape and as nothing else. The device
does not thereby compute a curvature, and this book does not claim it
does; whether that reading could ever be earned is a question that
stays on the far side of a fence this series has not crossed --- a face
still owed, like the fluid's. Illustration, not identity.
\end{fence}

Take the chapter's inventory, because it completes an arc. A machine
that measures the failure of loops to close; a second route to it
through the failure of probes to commute; a definition of flatness with
honest content; and all of it --- the reader should say this aloud,
because few textbooks pause on it --- without a single length. Nothing
so far prefers one scale to another, no angle has been measured, no
distance compared, and yet charts, arrows, machines, carrying, and
curvature all stand. The construction has run out of things it can
build without a ruler exactly at the moment the reader can feel what a
ruler would be for.

So the next chapter, at last, goes shopping. It buys the metric --- the
smooth choice of an inner product at every point --- and the reader will
watch what the purchase settles all at once: lengths and angles, the
marriage of arrows to read-outs promised in Chapter 3, and, by the
fundamental theorem promised in Chapter 5, the selection of exactly one
carrying rule from among all the candidates. The ruler arrives last
because everything before it never needed one; what it buys is the
subject of the next chapter, and the price is paid in the open.
